% Part: computability
% Chapter: recursive-functions
% Section: non-pr-functions

\documentclass[../../../include/open-logic-section]{subfiles}

\begin{document}

\olfileid{cmp}{rec}{npr}
\olsection{আদিম পুনরাবৃত্ত নয় এমন অপেক্ষক}

আদিম পুনরাবৃত্ত অপেক্ষকগুলি স্বজ্ঞাতভাবে গণনসাধ্য সব অপেক্ষককে নিঃশেষ
করে না। সব এক-স্থানীয় আদিম পুনরাবৃত্ত অপেক্ষকের এমন একটি তালিকা
$f_0$, $f_1$, $f_2$,~\dots যে বানানো যায়, যাতে নিবেশ $y$-এ $f_x$-এর
মান আমরা কার্যকরভাবে গণনা করতে পারি, তা স্বজ্ঞাতভাবে স্পষ্ট হওয়া উচিত;
অন্যভাবে বললে, নিচের সমীকরণে সংজ্ঞায়িত অপেক্ষক $g(x,y)$ গণনসাধ্য:
\[
g(x,y) = f_x(y)
\]
কিন্তু তাহলে নিচের অপেক্ষকটিও গণনসাধ্য:
\begin{eqnarray*}
h(x) & = & g(x,x) + 1 \\
& = & f_x(x) +1.
\end{eqnarray*}
প্রত্যেক আদিম পুনরাবৃত্ত অপেক্ষক $f_i$-এর ক্ষেত্রে $i$-তে $h$ ও
$f_i$-এর মান আলাদা। অতএব $h$ গণনসাধ্য, কিন্তু আদিম পুনরাবৃত্ত নয়;
$g$-এর ক্ষেত্রেও একই কথা বলা যায়। এটি কান্টরের কর্ণীকরণ যুক্তির একটি
“কার্যকর” রূপ।

আদিম পুনরাবৃত্ত নয় এমন গণনসাধ্য অপেক্ষকের আরও প্রত্যক্ষ উদাহরণ দেওয়া
যায়। যেমন, $g^n(x)$ সংকেতলিপি দিয়ে মোট $n$টি $g$-সহ
$g(g(\dots g(x)))$ বোঝানো যাক; এবং অপেক্ষকগুলির একটি অনুক্রম
$g_0,g_1,\dots$ এভাবে সংজ্ঞায়িত করা যাক:
\begin{eqnarray*}
g_0(x) & = & x+1 \\
g_{n + 1}(x) & = & g_n^x(x)
\end{eqnarray*}
প্রত্যেক অপেক্ষক $g_n$ যে আদিম পুনরাবৃত্ত, তা যাচাই করা যায়। প্রতিটি
পরবর্তী অপেক্ষক তার আগেরটির চেয়ে অনেক দ্রুত বাড়ে; $g_1(x)$ হলো $2x$,
$g_2(x)$ হলো $2^x \cdot x$, এবং $g_3(x)$ মোটামুটি $x$টি $2$-এর ঘাতস্তূপের
মতো দ্রুত বাড়ে। আকারমান--পেতের অপেক্ষকটি মূলত $G(x) = g_x(x)$ অপেক্ষক,
এবং দেখানো যায় যে এটি যেকোনো আদিম পুনরাবৃত্ত অপেক্ষকের চেয়ে দ্রুত বাড়ে।

আদিম পুনরাবৃত্ত অপেক্ষকগুলির তালিকায়নের প্রসঙ্গে ফেরা যাক। মনে করুন,
প্রত্যেক আদিম পুনরাবৃত্ত অপেক্ষককে আমরা একটি প্রতীকী সংকেতলিপি দিয়েছি;
তাই সংকেতলিপিগুলিকে তালিকায়িত করাই যথেষ্ট। প্রত্যেক সংকেতলিপি $F$-কে
নিচের মতো পুনরাবৃত্তভাবে একটি স্বাভাবিক সংখ্যা $\#(F)$ দেওয়া যায়:
\begin{eqnarray*}
\#(0) & = & \langle 0 \rangle \\
\#(S) & = & \langle 1 \rangle \\
\#(\Proj{n}{i}) & = & \langle 2, n, i \rangle \\
\#(\fn{Comp}_{k,l}[H,G_0,\dots,G_{k-1}]) & = & \langle
3,k,l,\#(H),\#(G_0),\dots,\#(G_{k-1}) \rangle \\
\#(\fn{Rec}_l[G,H]) & = & \langle 4, l, \#(G), \#(H) \rangle
\end{eqnarray*}
এখানে আগের অংশের সংকেত ব্যবহার করে সংখ্যার প্রত্যেক অনুক্রমকে স্বাভাবিক
সংখ্যা হিসেবে দেখা যায়—এই তথ্যটি আমরা ব্যবহার করছি। এর ফল হলো, প্রত্যেক
সংকেতকে একটি স্বাভাবিক সংখ্যা দেওয়া হয়। অবশ্য কিছু অনুক্রম, এবং তাই কিছু
সংখ্যা, কোনো সংকেতলিপির সঙ্গে সংশ্লিষ্ট নয়; কিন্তু $i$ দিয়ে সংকেতায়িত
কোনো সংকেতলিপি থাকলে, অর্থাৎ $i$ যদি এমন কোনো সংকেতলিপির সংকেত হয়, তবে
$f_i$-কে সেই সংকেতলিপির এক-স্থানীয় আদিম
পুনরাবৃত্ত অপেক্ষক এবং অন্যথায় ধ্রুবক $0$ অপেক্ষক ধরা যায়। শেষ পর্যন্ত
এতে এক-স্থানীয় আদিম পুনরাবৃত্ত অপেক্ষকগুলিকে তালিকায়িত করার একটি
প্রত্যক্ষ উপায় পাওয়া যায়।

(আসলে ধ্রুবক শূন্য অপেক্ষকের মতো কিছু অপেক্ষক তালিকায় একাধিকবার আসবে।
এটি কেবল আমাদের সংকেতায়নের কৃত্রিম ফল নয়; ধ্রুবক শূন্য অপেক্ষকের
একাধিক সংকেতলিপি থাকারও ফল। পরে আমরা দেখব যে গণনসাধ্য উপায়ে এই
পুনরাবৃত্তিগুলি এড়ানো যায় না; যেমন, প্রদত্ত কোনো সংকেতলিপি ধ্রুবক শূন্য
অপেক্ষককে প্রকাশ করে কি না, তা নির্ণয়কারী কোনো গণনসাধ্য অপেক্ষক নেই।)

এখন আমরা $g(x,y)$-কে $f_x(y)$ দিয়ে প্রদত্ত অপেক্ষক হিসেবে নিতে পারি,
যেখানে $f_x$ বলতে সদ্য বর্ণিত তালিকায়নটি বোঝায়। $g(x,y)$ যে গণনসাধ্য,
তা আমরা কীভাবে জানি? স্বজ্ঞাতভাবে এটি স্পষ্ট: $g(x,y)$ গণনা করতে প্রথমে
$x$-এর সংকেত “ভেঙে” দেখে নিতে হবে সেটি কোনো এক-স্থানীয় অপেক্ষকের
সংকেতলিপি কি না। হলে, নিবেশ~$y$-এ সেই অপেক্ষকের মান গণনা করতে হবে।

\begin{tagblock}{TMs}
\begin{digress}
আপনি হয়তো ইতিমধ্যেই নিশ্চিত যে কিছুটা পরিশ্রম করে এমন একটি প্রোগ্রাম
(যেমন Java বা C++-এ) লেখা যায়; এবার আমরা চার্চ--টুরিং থিসিসের আশ্রয়
নিতে পারি, যা বলে স্বজ্ঞাতভাবে গণনসাধ্য যেকোনো কিছু একটি টুরিং যন্ত্র দিয়ে
গণনা করা যায়।

অবশ্য $g(x,y)$ যে গণনসাধ্য তা দেখানোর আরও প্রত্যক্ষ উপায় হলো একে
গণনাকারী একটি টুরিং যন্ত্র স্পষ্টভাবে বর্ণনা করা। বিশেষত এতে
চার্চ--টুরিং থিসিস ও স্বজ্ঞার আশ্রয় এড়ানো যাবে। শিগগিরই আমরা এতটা
উপকরণ তৈরি করব যে টুরিং যন্ত্রে \emph{অনুকরণ} করা যায় এমন একটি
গণনামূলক মডেল—অর্থাৎ পুনরাবৃত্ত অপেক্ষক—ব্যবহার করে $g(x,y)$-কে গণনসাধ্য
দেখাতে পারব।
\end{digress}
\end{tagblock}

\end{document}
